\documentclass[10pt]{article}
\usepackage[margin=0.7in]{geometry}
\usepackage[T1]{fontenc}
\usepackage[utf8]{inputenc}
\usepackage{lmodern}
\usepackage{amsmath,amssymb,microtype}
\usepackage[hidelinks]{hyperref}
\setlength{\parindent}{0pt}
\setlength{\parskip}{0.3em}
\setlength{\abovedisplayskip}{5pt}
\setlength{\belowdisplayskip}{5pt}
\pagestyle{empty}
\begin{document}
\begin{center}\large\bfseries Updated Climate-Zone Heat Model\end{center}

\textbf{Full current model.}
Let $i$ index ZIP codes, $d$ days, $k(i)$ SES clusters, and $\mathrm{cz}[i]$
climate zones. Retain the existing standardized heat $H_{id}$, ZIP mean
$\bar H_i$, acute heat $x_{id}=H_{id}-\bar H_i$, and population exposure $n_{id}$.
The updated model for ED counts is
\begin{align}
Y_{id}\mid\theta&\sim\operatorname{NB}(n_{id}\mu_{id},r),\qquad r>0,\nonumber\\
\log\mu_{id}&=\alpha+\gamma_{k(i)}+\beta_{\mathrm{chronic}}\bar H_i
+\beta_{\mathrm{acute},k(i)}(H_{id}-\bar H_i)
+\beta_{m(d)}+\beta_{w(d)}\nonumber\\
&\quad+f_{\mathrm{doy}}(d)
+f_{\mathrm{acute},\mathrm{cz}[i]}(H_{id}-\bar H_i)+u_i,\label{eq:model}\\
f_{\mathrm{acute},\mathrm{cz}[i]}(\cdot)
&=f_{\mathrm{shared}}(\cdot)+f_{\mathrm{dev},\mathrm{cz}[i]}(\cdot),\nonumber\\
\beta_{\mathrm{acute},k}&=\beta_{\mathrm{acute}}+\delta_k,\qquad
u_i\mid\tau_{\mathrm{ZIP}}\overset{\mathrm{iid}}{\sim}N(0,\tau_{\mathrm{ZIP}}^{-1}).\nonumber
\end{align}
Here, $u_i$ is the ZIP-specific random intercept, allowing each ZIP to have
its own underlying ED rate after accounting for the other predictors.
The precision $\tau_{\mathrm{ZIP}}$ is its inverse variance; iid denotes
independent, identically distributed intercepts conditional on that precision.
The terms $m(d)$ and $w(d)$ denote month and weekday; $f_{\mathrm{doy}}$ is the
seasonal RW2 effect. Reference-cluster constraints remain
$\gamma_{k_0}=\delta_{k_0}=0$, with $k_0=6$.
The SES-specific heat slopes are linear; the RW2 heat curves supply the
nonlinear response.

\textbf{Negative-binomial likelihood and ZIP intercept.}
The negative-binomial likelihood allows daily ED counts to vary more than
a Poisson model would allow around the same predicted mean.
Its conditional mean is $M_{id}=n_{id}\mu_{id}$ and its conditional variance
is $M_{id}+M_{id}^2/r$. The size parameter $r$ controls this extra variation;
as $r\to\infty$, the likelihood approaches Poisson.

\textbf{Correlated climate-zone deviations.}
Independent replicated climate-zone RW2 deviations are replaced by grouped
deviations using \texttt{group = climate\_zone\_group} and
\texttt{control.group = list(model = "exchangeable")}.
Each of the five zones (20, 22, 28, 29 and 30) has an adjustment curve
$f_{\mathrm{dev},\mathrm{cz}}$ to the shared heat curve. These curves retain
a common curvature precision $\tau_{\mathrm{dev}}$, with between-zone
correlation matrix
\[
R(\rho)=
\begin{pmatrix}
1&\rho&\rho&\rho&\rho\\
\rho&1&\rho&\rho&\rho\\
\rho&\rho&1&\rho&\rho\\
\rho&\rho&\rho&1&\rho\\
\rho&\rho&\rho&\rho&1
\end{pmatrix},
\qquad -\tfrac14<\rho<1.
\]
Each row and column represents one zone. Every pair shares correlation
$\rho$ between its adjustment curves, after accounting for the shared
curve. The bounds ensure a valid, nonsingular five-zone correlation matrix.
For ZIPs without a zone assignment, the zone deviation is zero; the other
predictor terms remain.

\textbf{PC hyperpriors.}
The seasonal RW2, shared-heat RW2, zone-deviation RW2 and ZIP-intercept
precisions each have a PC prior. For component $j$, write its scale as
$\sigma_j=\tau_j^{-1/2}=1/\sqrt{\tau_j}$; the baseline calibration is
\[
\Pr(\sigma_j>1)=0.01,\qquad
j\in\{\mathrm{doy},\mathrm{shared},\mathrm{dev},\mathrm{ZIP}\}.
\]
Here $j$ labels the four components; each has its own estimated precision.
For ZIP intercepts, $\sigma_j$ is their SD. For scaled RW2 terms, it controls
the penalized curve variation: these priors discourage bending, while
straight-line trends remain unpenalized by the RW2 prior.
The correlation prior favors $\rho=0$, with prior probability $0.5$ on
negative correlations and $0.5$ on positive correlations. Conditional on
either sign, the distance $d_\rho=\sqrt{-\log\det R(\rho)}$ from zero
correlation has an exponential distribution, calibrated so that
$\Pr(\rho>0.5)=0.1$. This prior is converted to INLA's computation scale
without changing its probabilities.
NB size $r$ uses \texttt{pc.mgamma} with INLA's default rate parameter 7,
shrinking toward the Poisson limit. The 39 non-intercept fixed coefficients
retain independent $N(0,1000)$ priors, where 1000 is the variance;
the overall intercept $\alpha$ remains flat.

\newpage

\textbf{Complete heat contrasts and excess rates.}
Derived effects include both the SES-specific linear slope and the two
heat-curve components. Compare acute heat $x=0$ (the ZIP's mean heat) with
$x=1$ (one county-wide heat SD above that mean):
\[
L_i=\beta_{\mathrm{acute},k(i)}
+f_{\mathrm{acute},\mathrm{cz}[i]}(1)
-f_{\mathrm{acute},\mathrm{cz}[i]}(0),\qquad
\mathrm{IRR}_{1\mathrm{SD},i}=e^{L_i},\qquad
p_i=\Pr(L_i>0\mid Y).
\]
INLA linear-combination marginals yield medians, 95\% credible intervals and
probabilities for the 42 observed SES-cluster--zone combinations, including missing zones.
Curve evaluations use the nearest nodes on the retained 50-node grid.

Excess-rate calculations also include the complete heat response, with the
SES-cluster monthly mean as reference, equal ZIP-day weights, and signed
differences per 10,000 population per day. The ZIP intercept and other
non-heat terms are held fixed within each contrast. Excess-rate uncertainty
was estimated using 2,000 joint posterior draws that respected the fitted
model constraints. Conditional sampling variances were checked against
those saved by INLA.

\textbf{Likelihood comparison and zero diagnostics.}
A matched Poisson model used the same 54,910 ZIP-days, population exposures,
predictor terms and common priors. DIC and WAIC were 65,207.33 and 65,254.86
for NB, compared with 91,394.49 and 94,721.52 for Poisson, supporting retention
of the NB baseline. Observed zeros numbered 47,898 (87.23\%); predicted zero
fractions were 86.98\% for NB and 74.16\% for Poisson. Zero probabilities
were evaluated at posterior mean fitted rates and, for NB, posterior mean
size. These are in-sample plug-in diagnostics, not fully integrated
posterior-predictive tests; checks also covered SES clusters, ZIPs, zones
and months.

\textbf{Precision-prior sensitivity.}
The baseline is retained as the primary model. The four precision priors
are varied jointly: stronger PC priors use
$\Pr(\tau_j^{-1/2}>0.5)=0.01$, weaker PC priors use
$\Pr(\tau_j^{-1/2}>2)=0.01$, and a comparison uses
$\tau_j\sim\operatorname{Gamma}(1,5\times10^{-5})$ with shape--rate
parameters. Correlation, NB-size and fixed-coefficient priors remain at
baseline. These fit-only comparisons use heat-effect marginals,
conditional SES-cluster IRRs, DIC and WAIC, without using additional
excess-rate draws.

\textbf{Coefficient-variance sensitivity.}
Separate NB fits changed only the Normal prior variance for the 39
non-intercept fixed coefficients from 1000 to 100 or 10, retaining zero prior
means, the flat intercept and all PC priors. Variance 100 was optimized
again from its fitted solution; this retained fit converged without
likelihood-recovery warnings. For variances 100 and 10, respectively,
DIC was 65,207.28 and 65,206.02, and WAIC was 65,254.33 and 65,253.74.
Relative to baseline, the largest heat-IRR median changes were 0.77\% and
0.52\%, and the largest exceedance-probability changes were 1.14 and 0.81
percentage points. The retained variance-100 fit agreed with baseline on
whether each of the 42 heat intervals excluded 1. At variance 10, one borderline
interval (SES cluster 14, zone 29) changed from including to excluding 1.
Both alternatives predicted about 86.98--86.99\% zeros. These fit-only
checks used INLA marginals without additional excess-rate draws; the
variance-1000 NB baseline was retained.
\end{document}
